Published values are rarely commensurable, so side-by-side display is inappropriate;
Table~\ref{tab:famreq} states what each family requires and how far this corpus has taken it; the counts of comparable evaluations are below. Decoding protocols
overlap most, and only within a modality and a shared public label schema; probes otherwise report AUROC
on different concept sets, splits, label definitions, decoder capacities and controls, obscuring
differences of a few points. Discovery overlaps least: width, sparsity, layer and
interpretation vary, cross-seed stability is seldom reported and nominal concepts cannot be matched.
Intervention records permit individual inspection, but differing end-points, magnitudes and controls preclude
effect-size comparison. The most compatible stratum, six chest-radiograph studies using linear decoders on CheXpert or MIMIC-CXR, still differs in targets, label versions, splits and
regularisation, so it shows six distinct estimands rather than a poolable stratum, which the reporting
checklists of \S\ref{sec:future} would supply. Independent replication by a second group is uncoded and
apparently rare, and same-data medical--general comparisons confound domain with data, objective,
architecture and protocol (\S\ref{sec:findings}).

\textbf{Modality and family.} Pathology appears in \Npathstud{} of the \Ncore{} studies and chest radiography in
\Ncxrstud{} (\Npathrec{} and \Ncxrrec{} records); every other modality is markedly rarer
(Fig.~\ref{fig:corpus}c). Chest radiography supplies \Nintcxr{} of the \Nintrec{} intervention records and
\Ntracecxr{} of the \Ntracerec{} tracing records, so multimodal-LLM interpretability is largely a
chest-radiograph literature; pathology dominates geometry and discovery, entering intervention only through
frozen-embedding post-processing. Modality is retained in the evidence tables because a claim rarely
transfers across modalities.

\textbf{Study quality.} Study-quality coding comes from full texts (four from abstracts,
\S\ref{sec:method}). The fields are the subset of the imaging-AI and clinical-AI reporting
checklists~\cite{mongan2020claim,norgeot2020minimum} that a representation-level claim can be scored on from a paper: they
record what was reported, not whether it was done well.
Of the \Ncore{} core studies, \Npreprint{} are preprints, \Nconf{} conference
papers, \Nwork{} workshop papers and \Njour{} journal articles; \Ngrey{} of the preprints are grey-literature deposits and
\Nsingle{} are single-author, admitted under the rule of \S\ref{sec:method} that any retrievable full text meeting the eligibility criteria is coded. Excluding the
grey deposits lowers the corpus-breadth tier of F8a$''$ alone and changes no direction; excluding all \Ngreysingle{} lowers F3c as well. Of these studies, \Nextyes{} evaluate on a cohort or
dataset unused for probe, dictionary or intervention training, \Nextno{} do not and \Nextunk{} are
indeterminate;
patient-, slide- or subject-level splits are stated in \Nsplityes{} and unstated in \Nsplitno{}. Seeds,
repeated runs, intervals or tests accompany the main result in \Nseedyes{}, not in \Nseedno{}, and
indeterminately in \Nseedunk{}. Of the \Nclaims{} records, \Ndirpos{} report a positive result,
\Ndirmix{} a mixed or partial one and \Ndirnull{} a null or negative one; null results this rare warrant
publication- and reporting-bias caution for every statement below. The young corpus relies more on internal
than external validation, a second reason \S\ref{sec:findings} grades breadth rather than pooling effect
sizes.

\textbf{Concept provenance and concept validity.} Among those records, \Nprovlabel{} use structured
schemas, \Nprovtext{} text-mined concepts, \Nprovexpert{} expert vocabularies, \Nprovmodel{}
model-generated banks and \Nprovnone{} no prior concept, \Nposthocname{} of the last naming features post
hoc. Expert-defined and model-generated concepts occur almost only in intrinsic designs
(\S\ref{sec:intrinsic}). The best-supported decoding findings (\S\ref{sec:findings}, F1a and F2) rest on
labels automatically extracted from chest-radiograph
reports~\cite{irvin2019chexpert,johnson2019mimicjpg}, so concept validity reflects a labeller, not clinician
judgement (\S\ref{sec:eval-data}). Discovery without a prior concept shifts the burden to post-hoc naming,
clinician-validated by fewer than half of dictionary studies (\S\ref{sec:eval-what}). Provenance locates where a claim may fail for reasons unrelated to the model.
